\section{The remaining growing-host size as an elliptic-unit index sum}
\label{sec:elliptic-unit-budget}

\begin{theorem}[AC-K27: an exact numerical gate over the fixed base]
\label{AC-K27}
On each actual host $H=F_x(\sqrt[3]{269})$ of AC-K25,
use the derived elliptic-unit indices of AC-K26-elliptic
and put
\[
 \mathcal I_x=
 \sum_{1\leq e<t_x}\sum_\ell
       \frac{\log\#k_{O_\ell}}{b_\ell}
            \sum_{\eta\in\widehat G}\max(0,i_{\ell,\eta,e}),
       \qquad t_x=\lfloor x^{401}\rfloor .           \tag{EB.1}
\]
This is a finite nonnegative sum.  The all-prime
strict size of the growing field satisfies
\[
 \left|\sum_{1\leq e<t_x}\sum_\ell\log\#S_\ell(H,e)
                                      -\mathcal I_x\right|
    =O_\varepsilon\!\left(
             x^{806+16\varepsilon}(\log x)^2\right)
    =o\!\left(\exp(x^\varepsilon/2)\right).          \tag{EB.2}
\]
Consequently a little-o estimate at this exponential
scale for the growing-field strict sum is equivalent
to the same estimate for $\mathcal I_x$.
At a prescribed positive budget constant, either
estimate transfers with any fixed positive relative
slack, for sufficiently large $x$.

All finite Fourier, local invariant, local-condition,
coefficient-normalization and degree-one torsion costs
in this reduction have been bounded.  The unproved
size assertion is the estimate for the actual indices
in (EB.1); it is not replaced by determinant equality.
\end{theorem}
\begin{proof}
Let $\mathcal B_x$ be (FC.15), and define
\[
 \mathcal H_x=
 \sum_{1\leq e<t_x}\sum_\ell\frac1{b_\ell}
             \sum_{\eta\in\widehat G}
                    \log\#H^2(G_{K,\Sigma_\ell},T_{\eta,e}).
\]
It is finite by AC-K26-local.  Equations (EU.1)--(EU.2)
and AC-K26-torsion give
\[
 0\leq\mathcal H_x-\mathcal I_x
   \leq\sum_{1\leq e<t_x}\sum_\ell\frac1{b_\ell}
            \sum_\eta\log\#H^1(G_{K,\Sigma_\ell},T_{\eta,e})_{\rm tors}
   =O_\varepsilon\!\left(
              x^{806+16\varepsilon}(\log x)^2\right).
                                                            \tag{EB.3}
\]
In particular (EB.1) has finite support and no
cancellation of divergent or signed prime sums is used.
AC-K26-local gives the same error order between
$\mathcal H_x$ and $\mathcal B_x$.
AC-K25-budget compares $\mathcal B_x$ with the
growing-field strict sum, with error
$O_\varepsilon(x^{409+161\varepsilon/5}(\log x)^2)$.
For $0<\varepsilon<1/52$ this latter exponent is
strictly smaller than $806+16\varepsilon$.
The triangle inequality proves (EB.2).
Its little-o conclusion and both transfer assertions
follow by dividing by $\exp(x^\varepsilon/2)$.
\end{proof}

\paragraph{What is and is not closed.}
The geometric objects in this reduction are explicit:
a finite abelian cover of the arithmetic curve over
$K=\mathbb Q(\omega)$, its integral Fourier coefficient
sequence, and the inverse determinant line of a
rank-one Artin--Tate cohomology complex over its
two-variable Iwasawa weight space.  The elliptic-unit
theorem fixes the generator of that line.  The
valuation of its derived specialization relative to
the integral free cohomology lattice is still a
numerical quantity to estimate.

Equation (EB.2) closes the comparison to that quantity;
it does not close its upper bound, ALP\_MACHINE H04.4,
or the cubic Riemann hypothesis.  Nor does it turn
the source's one-way theta/congruence upper bound
into a Selmer upper bound.  Identification of the
prescribed mixed theta and Selmer data, and the
source's common-window quantifiers and positive
slack, remain separate requirements.  The SGA~1
and K-theory dictionaries already attached to
AC-K25 and AC-K10 apply to these objects; no new
Grothendieck archive scan or surface construction
is asserted here.
